\documentclass[11pt]{article}
\usepackage[margin=1in]{geometry}
\usepackage{amsmath,amssymb,amsthm,mathtools}
\usepackage[T1]{fontenc}
\usepackage{lmodern}
\usepackage{graphicx}
\usepackage{hyperref}
\usepackage{xcolor}

\newtheorem{theorem}{Theorem}
\newtheorem{corollary}[theorem]{Corollary}
\newtheorem{lemma}[theorem]{Lemma}
\newtheorem{remark}[theorem]{Remark}

\newcommand{\Om}{\Omega}
\newcommand{\R}{\mathbb R}
\newcommand{\Div}{\operatorname{div}}
\newcommand{\Breg}{\mathcal B_F}
\newcommand{\Erel}{\mathcal E_F}
\newcommand{\J}{\mathcal J_1}

\title{Relative-Bregman uniqueness for the $L^1$-gradient flow\\
and the full strongly elliptic quadratic subcase}
\author{Codex agent lane 8 (GPT5.6)}
\date{22 July 2026}

\begin{document}
\maketitle

\begin{center}
\fbox{\parbox{0.91\linewidth}{\textbf{Status.}
Candidate partial result, likely valid, human review requested.  The result
does not settle uniqueness for every weak minimizing-movement limit of a
general strongly convex integrand.}}
\end{center}

\section{Source question}

Chambolle and Novaga study the $L^1$-minimizing movement for
\[
  \Phi(u)=\int_\Om F(Du)
\]
on a bounded connected domain, with homogeneous Neumann or Dirichlet boundary
conditions.  They obtain limit solutions, characterize them for smooth
strongly convex $F$, prove uniqueness for monotone solutions, and prove
unconditional uniqueness for the isotropic Dirichlet energy
$F(\xi)=|\xi|^2/2$ \cite{CN}.  Their abstract states:

\begin{center}
\includegraphics[width=0.97\linewidth]{figures/open_problem_crop.png}
\end{center}

\begin{quote}
``Under a monotonicity assumption we can also prove uniqueness of the limit
solution, even though this remains an open question in full generality.''
\end{quote}

This packet proves two results toward that question.  First, a relative
Bregman estimate gives uniqueness for genuinely nonlinear uniformly convex
integrands whenever one of the two solutions has an integrable
$L^\infty_x$ bound on $D\dot v$.  Second, the nonlinear error vanishes
identically for every constant-coefficient strongly elliptic quadratic
integrand.  This yields unconditional uniqueness, contraction, and
whole-family minimizing-movement convergence for that complete anisotropic
quadratic class.

\section{Framework and definitions}

Let $\Om\subset\R^d$ be bounded, connected, and Lipschitz.  Put
\[
  \Psi(w)=\frac12\|w\|_{L^1(\Om)}^2,
  \qquad \J=\partial\Psi\subset L^1(\Om)\times L^\infty(\Om).
\]
For $w\ne0$, the relation $q\in\J(w)$ is equivalent to
\begin{equation}\label{eq:Jchar}
  \|q\|_\infty=\|w\|_1,
  \qquad qw=\|w\|_1|w|\quad\text{a.e.},
\end{equation}
while $\J(0)=\{0\}$.  In particular, $\J$ is monotone:
\begin{equation}\label{eq:Jmono}
  \int_\Om(q_1-q_2)(w_1-w_2)\,dx\ge0
  \quad(q_i\in\J(w_i)).
\end{equation}

Write $A=\nabla F$ and $H=D^2F$.  A sufficiently regular solution of the
flow characterized in \cite[Theorem 2.15]{CN} satisfies
\begin{equation}\label{eq:flow}
  q_u:=\Div A(Du)\in\J(\dot u)
  \quad\text{for a.e. time}.
\end{equation}
The pointwise formulation in the source paper,
\[
 |q_u|\le\|\dot u\|_1,
 \qquad \dot u\,q_u=|\dot u|\|\dot u\|_1,
\]
is precisely \eqref{eq:flow}.  We use either common Dirichlet data, so that
$u-v$ and $\dot u-\dot v$ have zero trace, or the natural Neumann condition
$A(Du)\cdot n=A(Dv)\cdot n=0$.

For $a,b\in\R^d$, define the one-sided Bregman divergence
\[
  \Breg(a\mid b)=F(a)-F(b)-A(b)\cdot(a-b)
\]
and the relative energy
\[
  \Erel(u\mid v)=\int_\Om \Breg(Du\mid Dv)\,dx.
\]

\section{Proof intuition}

Monotonicity of the $L^1$ duality map already gives the favorable term
\[
 \int_\Om [A(Du)-A(Dv)]\cdot D(\dot u-\dot v)\,dx\le0.
\]
For a quadratic $F$, this is exactly the time derivative of the energy of
$D(u-v)$, which is the contraction used in the source paper.  For nonlinear
$F$, differentiating the relative Bregman energy produces one additional
Taylor remainder:
\[
 [A(Du)-A(Dv)-H(Dv)(Du-Dv)]\cdot D\dot v.
\]
If $H$ is Lipschitz, this is quadratic in $Du-Dv$.  Uniform convexity says
that the Bregman energy is also quadratic in that difference.  Thus an
integrable $L^\infty$ bound on $D\dot v$ closes a Gronwall inequality.

This mechanism is asymmetric on purpose: it is a weak--strong estimate.
It also identifies the exact obstruction to using the quadratic proof in
the full weak class.  The regularity obtained in \cite{CN} controls
$D\dot v$ in $L^2$, not in $L^1_tL^\infty_x$, so the theorem below does not
silently resolve the remaining open case.

\section{The relative-Bregman estimate}

\begin{theorem}[Nonlinear weak--strong uniqueness]\label{thm:main}
Assume $F\in C^2(\R^d)$ and that there are constants
$0<\gamma\le\Gamma<\infty$ and $K<\infty$ such that
\begin{align}
  \gamma I&\le H(\xi)\le\Gamma I
       &&(\xi\in\R^d), \label{eq:elliptic}\\
  \|H(\xi)-H(\eta)\|&\le K|\xi-\eta|
       &&(\xi,\eta\in\R^d). \label{eq:hesslip}
\end{align}
Let $u,v\in AC([0,T];H^1(\Om))$ satisfy \eqref{eq:flow}, the same boundary
condition, and suppose
\[
  D\dot v\in L^1(0,T;L^\infty(\Om)).
\]
Assume the integrations by parts in \eqref{eq:pairing} below are valid
(for example, $q_u,q_v\in L^1(0,T;L^\infty)$ with the stated boundary
conditions).  Then for $0\le s\le t\le T$,
\begin{equation}\label{eq:gronwall}
  \Erel(u(t)\mid v(t))
  \le \Erel(u(s)\mid v(s))
     \exp\!\left(\frac K\gamma
       \int_s^t\|D\dot v(r)\|_\infty\,dr\right).
\end{equation}
Consequently, two such solutions with the same initial datum coincide.
\end{theorem}

\begin{proof}
Set $a=Du$ and $b=Dv$.  The hypotheses on $H$ imply the quadratic
equivalence
\begin{equation}\label{eq:bregcoerc}
  \frac\gamma2|a-b|^2\le\Breg(a\mid b)
  \le\frac\Gamma2|a-b|^2.
\end{equation}
The Banach-space chain rule (or a direct calculation for smooth curves,
followed by approximation) gives
\begin{align}
 \frac d{dt}\Breg(a\mid b)
 &= [A(a)-A(b)]\cdot(\dot a-\dot b) \notag\\
 &\quad+
 [A(a)-A(b)-H(b)(a-b)]\cdot\dot b. \label{eq:bregidentity}
\end{align}
After integration over $\Om$, the first term is
\begin{align}
 P(t)&:=\int_\Om[A(Du)-A(Dv)]\cdot
                   D(\dot u-\dot v)\,dx \notag\\
     &=-\int_\Om(q_u-q_v)(\dot u-\dot v)\,dx\le0. \label{eq:pairing}
\end{align}
Here the equality is integration by parts, and the inequality is the
monotonicity \eqref{eq:Jmono}.

The integral form of Taylor's theorem and \eqref{eq:hesslip} yield
\begin{equation}\label{eq:taylor}
 \big|A(a)-A(b)-H(b)(a-b)\big|
 \le\frac K2|a-b|^2.
\end{equation}
Combining \eqref{eq:bregidentity}--\eqref{eq:taylor} and using
\eqref{eq:bregcoerc}, we obtain for a.e. $t$,
\[
 \frac d{dt}\Erel(u\mid v)
 \le \frac K2\|D\dot v\|_\infty
       \int_\Om|Du-Dv|^2\,dx
 \le \frac K\gamma\|D\dot v\|_\infty\Erel(u\mid v).
\]
Gronwall's inequality proves \eqref{eq:gronwall}.  If $u(0)=v(0)$,
then \eqref{eq:bregcoerc} gives $Du(t)=Dv(t)$ for every $t$.

With Dirichlet data, $u=v$ follows immediately.  With Neumann data,
$u(t)-v(t)$ is spatially constant.  Lemma~\ref{lem:constant} below shows
that its time derivative is zero; the initial equality then gives $u=v$.
\end{proof}

\begin{lemma}[Neumann constant rigidity]\label{lem:constant}
Let $\Om$ be connected.  If $f\in H^1(\Om)$, $c\in\R$, and
$q\in L^\infty(\Om)$ satisfy
\[
  \int_\Om q\,dx=0,
  \qquad q\in\J(f)\cap\J(f+c),
\]
then $c=0$.
\end{lemma}

\begin{proof}
If $q=0$, then \eqref{eq:Jchar} gives $f=f+c=0$.  Suppose $q\ne0$.
Since $q$ has mean zero, the sets $\{q>0\}$ and $\{q<0\}$ have positive
measure.  Relation \eqref{eq:Jchar} gives
\[
 f\ge0\ \text{on }\{q>0\},
 \qquad f+c\le0\ \text{on }\{q<0\}.
\]
If $c>0$, these say that $f$ lies in $[0,\infty)$ on one positive-measure
set and in $(-\infty,-c]$ on another.  Moreover, on $\{q=0\}$,
\eqref{eq:Jchar} for the two nonzero vectors forces both $f=0$ and
$f+c=0$; hence this set is null if $c\ne0$.  A Lipschitz truncation of
$f$ mapping $(-\infty,-c]$ to $0$ and $[0,\infty)$ to $1$ is therefore
an $H^1$ function taking only the values $0$ and $1$ a.e., with both
level sets of positive measure.  Such an $H^1$ function has zero weak
gradient and must be constant on connected $\Om$, a contradiction.
The case $c<0$ is symmetric.
\end{proof}

\begin{corollary}[A genuinely nonlinear family]\label{cor:nonlinear}
For every $\varepsilon>0$, the integrand
\[
 F_\varepsilon(\xi)=\frac12|\xi|^2+
   \varepsilon\sum_{j=1}^d\log\cosh(\xi_j)
\]
satisfies the hypotheses of Theorem~\ref{thm:main}.  Thus its flow is
unique in the class of that theorem.
\end{corollary}

\begin{proof}
Its Hessian is diagonal with entries
$1+\varepsilon\operatorname{sech}^2(\xi_j)$.  Hence
$I\le H\le(1+\varepsilon)I$, and the derivative of the Hessian is globally
bounded.  Theorem~\ref{thm:main} applies.
\end{proof}

\section{The full strongly elliptic quadratic subcase}

\begin{corollary}[Anisotropic quadratic contraction and uniqueness]
\label{cor:quadratic}
Let $M$ be a constant symmetric positive-definite $d\times d$ matrix and
\[
 F(\xi)=\frac12\xi\cdot M\xi.
\]
For any two energy-class solutions of \eqref{eq:flow} with common
Dirichlet data or homogeneous Neumann data,
\begin{equation}\label{eq:quadraticcontraction}
 \frac12\int_\Om D(u(t)-v(t))\cdot M D(u(t)-v(t))\,dx
 \le
 \frac12\int_\Om D(u(s)-v(s))\cdot M D(u(s)-v(s))\,dx.
\end{equation}
No $L^\infty$ assumption on $D\dot v$ is needed.  In particular the
flow from a fixed initial datum is unique.

Moreover, the unique discrete minimizers of
\[
 \Phi(w)+\frac1{2\tau}\|w-w^{n-1}\|_1^2,
 \qquad \Phi(w)=\frac12\int_\Om Dw\cdot M Dw,
\]
converge, without choosing a subsequence, to this flow as $\tau\downarrow0$
in the same topologies as the quadratic minimizing movements in \cite{CN}.
\end{corollary}

\begin{proof}
Here $H\equiv M$, so $K=0$ and the Taylor remainder in
\eqref{eq:bregidentity} vanishes identically.  Therefore
\[
 \frac d{dt}\frac12\int_\Om D(u-v)\cdot M D(u-v)\,dx
 =-\int_\Om(q_u-q_v)(\dot u-\dot v)\,dx\le0,
\]
which proves \eqref{eq:quadraticcontraction}.  Dirichlet uniqueness is
immediate; Neumann uniqueness follows from Lemma~\ref{lem:constant}.

The existence, time-regularity, and compactness estimates in Section 4 of
\cite{CN} use only polarization of the quadratic form, ellipticity, and
the monotonicity of $\J$.  Replacing $|D\cdot|^2$ by
$D\cdot\,M D\cdot$ changes the estimates only by the two ellipticity
constants of $M$.  Thus every vanishing-step sequence has a subsequence
converging to an energy-class solution.  Uniqueness shows that all cluster
points agree.  Relative compactness then upgrades subsequential convergence
to convergence of the whole family.
\end{proof}

\section{Verification notes and limitations}

\begin{enumerate}
\item \textbf{Sign check.}  The source equation has
$q=\Div\nabla F(Du)\in\J(\dot u)$.  Hence monotonicity of $\J$ makes
\(-\int(q_u-q_v)(\dot u-\dot v)\) nonpositive, exactly as required.

\item \textbf{Algebra check.}  The script
\texttt{code/check\_bregman\_identity.py} numerically differentiates the
relative Bregman energy for the nonlinear integrand in
Corollary~\ref{cor:nonlinear} and compares it with the two terms in
\eqref{eq:bregidentity}.  It also tests the Taylor remainder bound on
10,000 seeded random inputs.  This is a regression/sign check, not a proof.

\item \textbf{What remains open.}  The source paper proves only
$\dot u\in L^2_{\mathrm{loc}}((0,\infty);H^1(\Om))$ for its general
strongly convex limit solutions.  That does not imply
$D\dot u\in L^1_tL^\infty_x$.  Theorem~\ref{thm:main} therefore gives a
complete classical/weak--strong uniqueness criterion but not unconditional
uniqueness for every source-paper limit solution.

\item \textbf{Neumann issue.}  Relative energy controls gradients only.
Lemma~\ref{lem:constant} is essential; it rules out a time-dependent spatial
constant by using the $L^1$ duality relation and connectedness.

\item \textbf{Novelty status.}  A bounded search on 22 July 2026 used the
exact title, arXiv:2302.12786, DOI 10.1137/22M1527556, the quoted open phrase,
and combinations of ``$L^1$ gradient flow'', ``strongly convex'', and
``uniqueness''.  Sources checked included arXiv, the 2024 SIAM journal page,
CVGMT, and general web/arXiv search.  The journal version still states the
full question as open, and no later paper explicitly closing it was found.
The relative-Bregman argument may overlap standard weak--strong relative
entropy techniques, so novelty confidence for the method is moderate; the
mathematical claim is presented as a candidate partial result, not as a
priority claim.
\end{enumerate}

\section{Human-review recommendation}

The packet is recommended for expert review as a likely valid partial result.
The key checks are: (i) the chain-rule regularity in the desired solution
class; (ii) the Neumann constant-rigidity lemma; and (iii) that the source
paper's Section 4 estimates transfer to a constant SPD matrix exactly as
claimed.  A useful next push would prove smoothing strong enough to place
every positive-time limit flow in $L^1_tW^{1,\infty}_x$, or replace the
$L^\infty$ factor in the commutator estimate by a norm already controlled by
the source paper.

\begin{thebibliography}{9}
\bibitem{CN}
A. Chambolle and M. Novaga,
\emph{$L^1$-gradient flow of convex functionals},
SIAM J. Math. Anal. \textbf{56} (2024), no.~5, 5747--5781;
arXiv:2302.12786; DOI: 10.1137/22M1527556.
\end{thebibliography}

\end{document}
